\documentclass[11pt,a4paper]{article}
\usepackage{geometry}
\geometry{margin=1in}
\usepackage{amsmath, amssymb}
\usepackage{graphicx}
\usepackage{booktabs}
\usepackage{hyperref}

\title{\textbf{Topological Ablation Metrics}\\ \Large PR-AUC and RT-RMSE Evaluation ($N=128$)}
\author{Antigravity AI Pipeline}
\date{August 2026}

\begin{document}
\maketitle

\begin{abstract}
While the Curvature-Penalized Likelihood provided formal information-theoretic evidence of topological superiority, translating the geometric expansion of the Granule Cell manifold into functional predictive metrics is crucial for evaluating biological isomorphism. In this report, we benchmark the Intact ECCM against the Lesioned ECCM and the baseline WSLS models using two macroscopic behavioral metrics: PR-AUC (Precision-Recall Area Under the Curve) for choice accuracy, and RT-RMSE (Root Mean Square Error of Response Times) for reaction time kinetics.
\end{abstract}

\section{Methodology}
For each of the $N=128$ participants, we generated predicted choice probabilities $P(Choice=1)$ and expected reaction times $E[RT]$ directly from the Hierarchical MAP parameters. 

\begin{itemize}
    \item \textbf{PR-AUC:} Precision-Recall AUC is utilized instead of standard ROC-AUC due to the heavy temporal imbalance inherent in shift-register predictions during non-stationary transitional phases. Higher is better.
    \item \textbf{RT-RMSE:} Computes the squared deviation between the empirical human reaction times and the expected bounds produced by the Gated Drift-Diffusion framework. Lower is better.
\end{itemize}

\section{Aggregate Metric Comparisons}

\begin{table}[h!]
    \centering
    \begin{tabular}{lcc}
        \toprule
        \textbf{Architecture} & \textbf{Mean PR-AUC ($\uparrow$)} & \textbf{Mean RT-RMSE ($\downarrow$)} \\
        \midrule
        WSLS (Baseline) & 0.7213 & 0.6483 \\ 
        Lesioned ECCM (Ablated) & 0.8083 & 0.6170 \\ 
        \textbf{Intact ECCM (Unified)} & \textbf{0.8286} & \textbf{0.6092} \\ 
        \bottomrule
    \end{tabular}
    \caption{Macroscopic behavioral prediction metrics evaluated across $N=128$ participants.}
\end{table}

\section{Statistical Significance (Paired $t$-tests)}

\subsection{Intact ECCM vs. WSLS Baseline}
\begin{itemize}
    \item \textbf{PR-AUC Difference:} $+0.1073$ ($t = 24.743$, $p < 2.2 \times 10^{-16}$)
    \item \textbf{RT-RMSE Difference:} $-0.0391$ ($t = -17.833$, $p < 2.2 \times 10^{-16}$)
\end{itemize}
\textbf{Conclusion:} The high-dimensional geometric model dominates the static heuristic rule, yielding a massive 10.7\% boost in choice prediction and dropping RT error significantly.

\subsection{Intact ECCM vs. Lesioned ECCM (The Ablation Test)}
\begin{itemize}
    \item \textbf{PR-AUC Difference:} $+0.0203$ ($t = 7.904$, $p = 5.617 \times 10^{-13}$)
    \item \textbf{RT-RMSE Difference:} $-0.0077$ ($t = -6.285$, $p = 2.413 \times 10^{-9}$)
\end{itemize}
\textbf{Conclusion:} The geometric expansion is functionally necessary. Even in macroscopic terms, ablating the Granule Cell manifold ($1024$ nodes) forces the linear readout to collapse dimensional separation, sacrificing $2\%$ in PR-AUC precision and compounding temporal errors in the response distributions ($p < 10^{-8}$). 

The Cerebellar topology is formally required to execute optimal high-dimensional choice prediction and neuromodulatory timing constraints.
\end{document}
